% !TEX TS-program = pdflatex
\documentclass[11pt,a4paper]{article}

\usepackage[margin=2.5cm]{geometry}
\usepackage[T1]{fontenc}
\usepackage{lmodern}
\usepackage{microtype}
\usepackage{setspace}
\onehalfspacing

\usepackage{amsmath,amssymb}
\usepackage{booktabs}
\usepackage{longtable}
\usepackage{enumitem}
\usepackage[hidelinks]{hyperref}
\usepackage{xurl}
% Allow line breaks in URLs at more places
\makeatletter
\g@addto@macro{\UrlBreaks}{\do\/\do-\do_\do.\do=\do?\do&}
\makeatother

\usepackage{titlesec}
\titleformat{\section}{\large\bfseries}{\thesection.}{0.5em}{}
\titleformat{\subsection}{\normalsize\bfseries}{\thesubsection}{0.5em}{}

\title{\textbf{Project Structure}\\[0.35em]
\large Atmospheric Water Harvesting for Autonomous PV Panels Cleaning ---\\
Methodology, Data, Tools, and Formulas}
\author{Graduation Project}
\date{\today}

\begin{document}
\maketitle
\tableofcontents
\newpage

\section{Purpose}
This document consolidates a structured workflow for the graduation project \textit{Atmospheric Water Harvesting for Autonomous PV Panels Cleaning}, coupling \textbf{solar PV performance (soiling)} with \textbf{water demand for autonomous cleaning} and \textbf{atmospheric water harvesting (AWH)} supply sizing. It lists authoritative links, recommended software, and core formulas with references suitable for a thesis bibliography. Those formulas are treated as the \textbf{foundation for selecting and building the preferred system configuration}, not as the final project deliverable by themselves.

\section{Methodology (Phased Workflow)}

\subsection{Phase A --- Reference plant class (not only a map pin)}
\begin{enumerate}[leftmargin=*]
  \item \textbf{Define plant classes} (state AC vs DC clearly):
  \begin{itemize}[nosep]
    \item Distributed / C\&I: typically $<5$--$10$\,MW.
    \item Medium utility: $\sim 10$--$50$\,MW (common in older FiT-era portfolios).
    \item Large utility: $50$--$500{+}$\,MW single-site or clustered mega-parks.
  \end{itemize}
  \item \textbf{Build a small inventory} (10--20 plants): MW, module technology (mono/bifacial), tracker/fixed tilt, ground coverage ratio (GCR), and qualitative dust proxies (roads, agriculture, quarries) from satellite imagery when possible.
  \item \textbf{Select a primary case study} matching the majority class you observe in Egypt's fleet; add a \textbf{sensitivity case} (e.g., windier coastal desert vs dustier Upper Egypt cluster).
\end{enumerate}
\textbf{Often overlooked constraints:} water rights and discharge rules, grid curtailment (affects revenue and surplus energy for AWH), O\&M access, and module warranty limits on cleaning methods.

\subsection{Phase B --- Weather and soiling (climate $\rightarrow$ cleaning need)}
\begin{enumerate}[resume,leftmargin=*]
  \item Obtain a \textbf{typical meteorological year (TMY)} at plant coordinates: irradiance, temperature, wind, relative humidity (RH), precipitation.
  \item Separate \textbf{soiling drivers}: low-rain months, wind direction vs dust sources, dew/fog cycles (particle adhesion), avian/agricultural proximity.
  \item Decide whether cleaning water demand is driven by \textbf{fixed schedule}, \textbf{threshold on predicted soiling loss}, or \textbf{hybrid}.
\end{enumerate}
\textbf{Often overlooked:} \textbf{rain as partial cleaning}; PV simulation tools treat soiling as factors you must justify.

\subsection{Phase C --- Cleaning process model $\rightarrow$ water demand}
\begin{enumerate}[resume,leftmargin=*]
  \item Fix a \textbf{cleaning technology path}: manual hose, truck-mounted demineralized spray, robotic micro-spray, brush+dry, etc.\ (sets wetted area, pressure, recovery, blowdown).
  \item Compute \textbf{liters per event} and \textbf{events per year}; sum annual demand with a \textbf{safety margin} (commonly 20--40\%) for flushing, line losses, and quality excursions.
\end{enumerate}

\subsection{Phase D --- AWH screening (climate + duty + energy)}
\begin{enumerate}[resume,leftmargin=*]
  \item Shortlist AWH families for \textbf{hot arid} climates (often low daytime RH; better RH at night/dawn):
  passive dew/fog collectors; active cooling condensation / dehumidification; sorption-based systems (hygroscopic materials, brines, desiccants).
  \item For each option, estimate \textbf{seasonal water productivity} and \textbf{specific energy consumption (SEC)} [kWh/kg]. Compare against the cleaning schedule's \textbf{peak week} demand, not annual average only.
\end{enumerate}
\textbf{Often overlooked:} storage and water chemistry (TDS/hardness $\rightarrow$ spotting), night-only scheduling tied to RH/temperature cycles.

\subsection{Phase E --- Integrated simulation and prototype plan}
\begin{enumerate}[resume,leftmargin=*]
  \item \textbf{PV electrical annual yield and loss budget:} PVsyst and/or NREL SAM.
  \item \textbf{AWH subsystem:} usually external to PVsyst; TRNSYS, Modelica, or Python coupled to the same TMY.
  \item \textbf{CAD/CFD} for manifolds, nozzles, ducts; scale prototypes using dimensionless groups (Re, NTU, mass-transfer analogies) where relevant.
  \item \textbf{Prototype KPI:} e.g.\ measured kg water per kWh at controlled inlet psychrometric state.
\end{enumerate}

\section{Recommended Egyptian Anchors (Case Studies)}
\begin{table}[h]
\centering
\begin{tabular}{@{}p{3.2cm}p{5.8cm}p{6.2cm}@{}}
\toprule
\textbf{Anchor} & \textbf{Rationale} & \textbf{Starting links} \\
\midrule
Benban Solar Park (Aswan Gov.) &
Clustered utility plots; strong arid soiling context; widely referenced. &
\url{https://en.wikipedia.org/wiki/Benban_Solar_Park}\\
&& \url{https://benbansolarpark.com/}\\
\midrule
Kom Ombo PV (Aswan area) &
Large single-site utility narrative; good for layout thinking. &
\url{https://acwapower.com/en/projects/kom-ombo-pv/}\\
\midrule
National / global statistics context &
Justify ``what Egypt mostly builds'' vs headline projects. &
\url{https://www.irena.org/publications}\\
\bottomrule
\end{tabular}
\end{table}

\section{Data Portals and Simulation Tools}
\begin{description}[style=nextline,leftmargin=0pt,labelwidth=!]
  \item[NASA POWER] \url{https://power.larc.nasa.gov/} --- Solar, temperature, humidity, wind, precipitation; useful for TMY-style student datasets.
  \item[Copernicus CDS (ERA5)] \url{https://cds.climate.copernicus.eu/} --- Reanalysis for hourly humidity, wind, precipitation on a grid cell.
  \item[Global Solar Atlas] \url{https://globalsolaratlas.info} --- Benchmark GHI/PVOUT comparisons across Egypt.
  \item[PVsyst documentation (soiling)] \url{https://www.pvsyst.com/help/project-design/array-and-system-losses/soiling-loss.html}
  \item[NREL SAM] \url{https://sam.nrel.gov/} --- Free techno-economic cross-check for PV.
  \item[NREA (Egypt)] \url{http://www.nrea.gov.eg/test/en/About/Intro} --- Institutional orientation; pair with project fact sheets for hard numbers.
\end{description}

\section{Formulas and References}

\subsection{Module area from plant data}
If module count $N$ and module length/width $L,W$ are known:
\begin{equation}
A_{\mathrm{pv}} = N\,L\,W .
\end{equation}
If only DC nameplate $P_{\mathrm{dc,STC}}$ and representative STC module efficiency $\eta_{\mathrm{mod}}$ are known:
\begin{equation}
A_{\mathrm{pv}} \approx \frac{P_{\mathrm{dc,STC}}}{\eta_{\mathrm{mod}}\,G_{\mathrm{STC}}},
\qquad G_{\mathrm{STC}} = 1000\ \mathrm{W/m^2}.
\end{equation}
\textbf{Reference (monitoring definitions):} IEC~61724-1 (PV system performance monitoring) --- obtain via \url{https://webstore.iec.ch/}.

\subsection{Annual cleaning water (intensity $\times$ schedule)}
Let $w$ be water intensity [L/m$^2$] per cleaning event (process-dependent), $N_{\mathrm{clean}}$ events/year, and $f_{\mathrm{margin}}$ fractional margin for flushing/spills:
\begin{equation}
V_{\mathrm{year}} = w\,A_{\mathrm{pv}}\,N_{\mathrm{clean}}\,\bigl(1+f_{\mathrm{margin}}\bigr).
\end{equation}
\textbf{Experimental process anchor (read the paper before quoting $w$):}\\
\url{https://www.sciencedirect.com/science/article/abs/pii/S0038092X20302450}

\subsection{Economic framing for cleaning timing (literature entry point)}
Compare cleaning cost to incremental revenue from restored production (optimization papers extend this). Starting reference:\\
\url{https://www.sciencedirect.com/science/article/abs/pii/S0306261921000039}

\subsection{Psychrometrics (moisture in air)}
Let $p$ be total pressure, $\phi$ relative humidity, $p_{ws}(T)$ saturation pressure at dry-bulb $T$. Water vapor partial pressure:
\begin{equation}
p_w = \phi\, p_{ws}(T).
\end{equation}
Humidity ratio (absolute humidity):
\begin{equation}
\omega = 0.622\,\frac{p_w}{p - p_w}
\quad [\mathrm{kg_{vapor}/kg_{dry\ air}}].
\end{equation}
\textbf{Reference:} ASHRAE Handbook---Fundamentals, Chapter~1 (Psychrometrics).\\
\url{https://www.ashrae.org/technical-resources/handbooks}

\subsection{AWH performance metrics}
Specific electric energy consumption:
\begin{equation}
\mathrm{SEC}_{\mathrm{el}} = \frac{E_{\mathrm{el}}}{m_w}
\quad [\mathrm{kWh/kg}].
\end{equation}
Water productivity (example form):
\begin{equation}
Y_w = \frac{m_w}{A_{\mathrm{AWH}}\,\Delta t}
\quad [\mathrm{kg/m^2/day}].
\end{equation}
\textbf{Critical metrics / comparability:} Kim et al., \textit{Energy Environ.\ Sci.}\ (2022) ---\\
\url{https://pubs.rsc.org/en/content/articlelanding/2022/ee/d2ee03079a}

\textbf{Broad technology review:}\\
\url{https://www.sciencedirect.com/science/article/abs/pii/S0360544223015803}

\section{Software Stack (Beyond PVsyst + CAD)}
\begin{itemize}[leftmargin=*]
  \item PV energy + losses: PVsyst; cross-check with SAM; optional Python pvlib sanity checks.
  \item AWH: TRNSYS / Modelica / custom Python coupled to TMY psychrometrics.
  \item CAD: SolidWorks / Fusion~360; CFD (optional): OpenFOAM / Ansys Fluent for distribution uniformity.
  \item GIS: QGIS for land-cover / road proximity as dust proxies.
\end{itemize}

\section{Closing Integration Statement}
Anchor the plant class with evidence, lock a coordinate-specific TMY, translate soiling into a justified cleaning schedule, compute annual liters with uncertainty, then size AWH against peak-week demand reporting SEC (and thermal energy if applicable). Clearly separate what PVsyst validates (energy/soiling sensitivity) from what you must validate experimentally or with a dedicated AWH model (mass transfer, water quality, O\&M).

\section{Selected Configuration and Traceable Build Path}
\subsection{Calculated best-option parameters}
\begin{itemize}[leftmargin=*]
  \item \textbf{Numerical best LCOW route:} absorption-style AWH at SEC $\approx 0.38$\,kWh/kg under the current literature-based assumptions.
  \item \textbf{Prototype-selected route:} modular \textbf{vapor-compression condensation} with polishing treatment, selected because it remains near the numerical optimum while being easier to procure and assemble.
  \item \textbf{Shared design basis:} 200\,MW DC Benban-class block, $\eta_{\mathrm{mod}}=0.215$, $w=1.8$\,L/m$^2$/event, $N_{\mathrm{clean}}=18$/yr, $f_{\mathrm{margin}}=0.25$, and annual water duty on the order of $3.77\times 10^4$\,m$^3$/yr.
  \item \textbf{Prototype target:} 200\,L/day condensate, 500\,L raw tank, 500\,L polished-water tank, and 2.0\,m$^3$/h delivery at 3.0\,bar for interface validation.
\end{itemize}

\subsection{Traceability matrix}
\begin{table}[h]
\centering
\small
\begin{tabular}{@{}p{4.1cm}p{4.0cm}p{6.6cm}@{}}
\toprule
\textbf{Source calculation} & \textbf{Decision} & \textbf{Traceable design output} \\
\midrule
Module-area equation & Fix plant cleaning footprint & Defines the cleaning duty that all AWH concepts must satisfy. \\
Annual water-demand equation & Size storage and production & Requires modular generation plus buffering instead of direct supply only. \\
SEC comparison & Rank candidate technologies & Keeps absorption as numeric benchmark and VCR as the preferred build route. \\
Water-quality analysis & Select treatment train & Fixes sediment filtration, carbon polishing, UV, and optional RO bypass. \\
Cleaning-interface requirements & Specify hydraulic discharge & Sets the prototype pump, manifold, and outlet instrumentation targets. \\
\bottomrule
\end{tabular}
\end{table}

\subsection{Simulation plan}
\begin{enumerate}[leftmargin=*]
  \item Run PVsyst with the same cleaning schedule used in the water-duty equations.
  \item Build an hourly AWH simulation from NASA POWER and ERA5 psychrometrics.
  \item Simulate compressor, fan, condensate, and storage behavior for summer, shoulder, and winter cases.
  \item Add sensitivity cases for low-RH nights, filter fouling, and missed cleaning events.
  \item Validate the design when condensate production, polished-water TDS, and cleaning-interface pressure all meet the target band.
\end{enumerate}

\subsection{Engineering drawings, BOM, and component specifications}
\begin{itemize}[leftmargin=*]
  \item \textbf{Drawings:} D-101 process flow, D-102 skid general arrangement, D-103 tank and drain layout, D-104 treatment manifold, D-105 electrical one-line, D-106 controls and instrumentation, D-107 cleaning-interface manifold.
  \item \textbf{BOM:} compressor package, evaporator coil, condenser, EC fans, 500\,L raw tank, 500\,L polished tank, 5\,$\mu$m filter, carbon stage, UV reactor, 2.0\,m$^3$/h pump, flow meter, TDS meter, T/RH sensors, level switches, PLC/datalogger, valves, and tubing.
  \item \textbf{Specifications:} food-safe or HDPE wetted surfaces, serviceable filter cartridges, laboratory electrical supply with interlocks, and maintenance access around the skid footprint.
\end{itemize}

\subsection{Build manual, safety, and quality control}
\begin{enumerate}[leftmargin=*]
  \item Assemble the refrigeration loop and complete a leak test before energization.
  \item Install treatment and storage hardware according to D-101 and D-104.
  \item Calibrate sensors and run dry then wet commissioning before full 24-hour tests.
  \item Apply lockout-tagout, refrigerant handling, electrical earthing, and sanitation protocols.
  \item Hold QC checkpoints at leak test, airflow verification, condensate-rate verification, TDS verification, and discharge-pressure verification.
\end{enumerate}

\subsection{Updated master project plan}
\begin{table}[h]
\centering
\small
\begin{tabular}{@{}p{2.0cm}p{4.1cm}p{3.0cm}p{2.0cm}p{3.6cm}@{}}
\toprule
\textbf{Milestone} & \textbf{Deliverable} & \textbf{Resources} & \textbf{Date} & \textbf{Success criteria} \\
\midrule
M1 & Freeze design basis and route ranking & PI + analyst & T0+2d & Inputs and traceability approved \\
M2 & Complete PVsyst and AWH simulations & PV analyst + controls student & T0+7d & Weather-linked performance model closed \\
M3 & Release drawing set D-101 to D-107 & CAD lead + mechanical lead & T0+10d & All interfaces dimensioned \\
M4 & Procure BOM and fabricate skid & Procurement + workshop technician & T0+17d & Components inspected and ready \\
M5 & Build and commission prototype & Mechanical + electrical + safety officer & T0+24d & Safe 24-hour run achieved \\
M6 & Validate against selected targets & Test engineer + QA lead & T0+28d & Production, quality, and pressure targets achieved \\
\bottomrule
\end{tabular}
\end{table}

\end{document}
